% Options for packages loaded elsewhere
\PassOptionsToPackage{unicode}{hyperref}
\PassOptionsToPackage{hyphens}{url}
\documentclass[
]{article}
\usepackage{xcolor}
\usepackage{amsmath,amssymb}
\setcounter{secnumdepth}{-\maxdimen} % remove section numbering
\usepackage{iftex}
\ifPDFTeX
  \usepackage[T1]{fontenc}
  \usepackage[utf8]{inputenc}
  \usepackage{textcomp} % provide euro and other symbols
\else % if luatex or xetex
  \usepackage{unicode-math} % this also loads fontspec
  \defaultfontfeatures{Scale=MatchLowercase}
  \defaultfontfeatures[\rmfamily]{Ligatures=TeX,Scale=1}
\fi
\usepackage{lmodern}
\ifPDFTeX\else
  % xetex/luatex font selection
\fi
% Use upquote if available, for straight quotes in verbatim environments
\IfFileExists{upquote.sty}{\usepackage{upquote}}{}
\IfFileExists{microtype.sty}{% use microtype if available
  \usepackage[]{microtype}
  \UseMicrotypeSet[protrusion]{basicmath} % disable protrusion for tt fonts
}{}
\makeatletter
\@ifundefined{KOMAClassName}{% if non-KOMA class
  \IfFileExists{parskip.sty}{%
    \usepackage{parskip}
  }{% else
    \setlength{\parindent}{0pt}
    \setlength{\parskip}{6pt plus 2pt minus 1pt}}
}{% if KOMA class
  \KOMAoptions{parskip=half}}
\makeatother
\usepackage{graphicx}
\makeatletter
\newsavebox\pandoc@box
\newcommand*\pandocbounded[1]{% scales image to fit in text height/width
  \sbox\pandoc@box{#1}%
  \Gscale@div\@tempa{\textheight}{\dimexpr\ht\pandoc@box+\dp\pandoc@box\relax}%
  \Gscale@div\@tempb{\linewidth}{\wd\pandoc@box}%
  \ifdim\@tempb\p@<\@tempa\p@\let\@tempa\@tempb\fi% select the smaller of both
  \ifdim\@tempa\p@<\p@\scalebox{\@tempa}{\usebox\pandoc@box}%
  \else\usebox{\pandoc@box}%
  \fi%
}
% Set default figure placement to htbp
\def\fps@figure{htbp}
\makeatother
\setlength{\emergencystretch}{3em} % prevent overfull lines
\providecommand{\tightlist}{%
  \setlength{\itemsep}{0pt}\setlength{\parskip}{0pt}}
\usepackage{bookmark}
\IfFileExists{xurl.sty}{\usepackage{xurl}}{} % add URL line breaks if available
\urlstyle{same}
\hypersetup{
  pdftitle={Mathematical Analysis Notes: Tutorial 1 and Later},
  hidelinks,
  pdfcreator={LaTeX via pandoc}}

\title{Mathematical Analysis Notes: Tutorial 1 and Later}
\author{}
\date{}

\begin{document}
\maketitle

\section{Tutorial 1}\label{tutorial-1}

\subsection{1 Proof, logic, and limit language}\label{proof-logic-and-limit-language}

\subsubsection{1.1 Basic terminology}\label{basic-terminology}

\begin{quote}
• An axiom is a statement, rule, or principle that is accepted as true without needing proof. • A theorem is an important statement established by proof; a less central result is often called a proposition or lemma. • A proof is a finite, logically valid argument using definitions, axioms, and results already proved. • A statement of the form \(P ⇒ Q\) is a conditional statement: \(P\) is the hypothesis and \(Q\) is the conclusion.
\end{quote}

\subsubsection{\texorpdfstring{1.2 The \(\varepsilon\)-\(\delta\) definition}{1.2 The \textbackslash varepsilon-\textbackslash delta definition}}\label{the-varepsilon-delta-definition}

\begin{quote}
Here \(\delta\) controls how close the input \(x\) is to \(x_0\), while \(\epsilon\) prescribes how close \(f(x)\) must be to \(ℓ\).
\end{quote}

A standard proof has the following shape:

1. Let \(\varepsilon>0\) be arbitrary.
2. Choose an explicit \(\delta>0\), usually by first estimating \(|f(x)-\ell|\) in terms of \(|x-x_0|\).
3. Assume \(0<|x-x_0|<\delta\) and verify \(|f(x)-\ell|<\varepsilon\).

\subsubsection{1.3 Direct proof, contradiction, and contraposition}\label{direct-proof-contradiction-and-contraposition}

\paragraph{Direct proof}\label{direct-proof}

To prove a statement of the form

\[
P\Rightarrow Q,
\]

assume that \(P\) is true and use definitions, axioms, and previously proved results until \(Q\) follows. This is the default method for most \(\varepsilon\)-\(\delta\) arguments.

\paragraph{Proof by contradiction}\label{proof-by-contradiction}

To prove a statement \(S\):

\begin{enumerate}
\def\labelenumi{\arabic{enumi}.}
\tightlist
\item
  Assume its negation, \(\neg S\).
\item
  Derive a contradiction, such as \(R\land\neg R\), an impossible inequality, or a conflict with a known theorem.
\item
  Conclude that \(\neg S\) is false, so \(S\) is true.
\end{enumerate}

For a uniqueness statement, assume that two distinct objects both satisfy the defining property. For a nonexistence statement, assume that the object exists and use its required properties to obtain a contradiction.

\paragraph{\texorpdfstring{Example: irrationality of \(\sqrt2\)}{Example: irrationality of \textbackslash sqrt2}}\label{example-irrationality-of-sqrt2}

Suppose, for contradiction, that

\[
\sqrt2=\frac{p}{q},
\]

where \(p,q\) are positive integers with no common factor. Squaring gives

\[
p^2=2q^2.
\]

Thus \(p^2\) is even, so \(p\) is even. Write \(p=2r\). Substitution gives

\[
4r^2=2q^2
\quad\Longrightarrow\quad
q^2=2r^2.
\]

Therefore \(q\) is also even. Hence \(p\) and \(q\) have the common factor \(2\), contradicting the assumption that \(p/q\) is in lowest terms. Consequently,

\[
\sqrt2\notin\mathbb{Q}.
\]

\paragraph{Proof by contraposition}\label{proof-by-contraposition}

To prove

\[
P\Rightarrow Q,
\]

it is enough to prove the logically equivalent statement

\[
\neg Q\Rightarrow\neg P.
\]

This method is especially useful when the negation of the conclusion gives more concrete information than the conclusion itself.

\subsubsection{1.4 Mathematical induction}\label{mathematical-induction}

Let \(P(n)\) be a statement for \(n\in\mathbb{Z}_{>0}\). The principle of mathematical induction says that if:

\begin{enumerate}
\def\labelenumi{\arabic{enumi}.}
\tightlist
\item
  \textbf{Base case:} \(P(1)\) is true;
\item
  \textbf{Inductive step:} for every positive integer \(n\), assuming \(P(n)\) is true implies that \(P(n+1)\) is true;
\end{enumerate}

then \(P(n)\) is true for every positive integer \(n\).

In the inductive step, clearly state the \textbf{induction hypothesis} \(P(n)\) before proving \(P(n+1)\). Do not assume the statement that is being proved.

\paragraph{\texorpdfstring{Example: sum of the first \(n\) positive integers}{Example: sum of the first n positive integers}}\label{example-sum-of-the-first-n-positive-integers}

We prove that

\[
1+2+\cdots+n=\frac{n(n+1)}{2}
\qquad(n\geq1).
\]

For \(n=1\), both sides equal \(1\). Now assume the formula holds for \(n=k\):

\[
1+2+\cdots+k=\frac{k(k+1)}{2}.
\]

Then

\[
\begin{aligned}
1+2+\cdots+k+(k+1)
&=\frac{k(k+1)}{2}+(k+1)\\
&=\frac{(k+1)(k+2)}{2}.
\end{aligned}
\]

Thus, the statement for \(k\) implies the statement for \(k+1\), so the formula holds for every \(n\geq1\).

\paragraph{Strong induction}\label{strong-induction}

The principle of strong induction replaces the induction hypothesis \(P(n)\) by

\[
P(1),P(2),\ldots,P(n)\quad\Longrightarrow\quad P(n+1).
\]

It is useful when the case \(n+1\) depends on several earlier cases rather than only on the immediately preceding case.

Ordinary induction, strong induction, and the well-ordering principle are equivalent principles. The well-ordering principle states that every nonempty subset of \(\mathbb{Z}_{>0}\) has a least element.

\subsection{Lecture 5}\label{lecture-5}

2026 9.16

\subsubsection{}\label{section-1}

\subsection{Exercise 1.1.26}\label{exercise-1.1.26}

Any polynomial of odd degree with real coefficients has at least one real root.

\subsubsection{Proof}\label{proof}

Let

\[
P(x)=\sum_{k=0}^{n}a_kx^k
\]

be a polynomial of degree \(n\) with real coefficients, where \(n\) is an odd integer and \(a_n\neq0\). Without loss of generality, assume that \(a_n>0\). If \(a_n<0\), we can apply the same argument to \(-P\).

Let

\[
C=\max_{0\leq k\leq n-1}|a_k|.
\]

For \(|x|\geq1\), we have

\[
\begin{aligned}
\left|P(x)-a_nx^n\right|
&=\left|\sum_{k=0}^{n-1}a_kx^k\right|\\
&\leq C\sum_{k=0}^{n-1}|x|^k\\
&\leq Cn|x|^{n-1}
=\frac{Cn}{|a_n||x|}|a_nx^n|.
\end{aligned}
\]

Choose

\[
R>\max\left\{1,\frac{2Cn}{|a_n|}\right\}.
\]

Then, whenever \(|x|\geq R\),

\[
\left|P(x)-a_nx^n\right|<\frac12|a_nx^n|.
\]

Since \(n\) is odd, \(R^n>0\) and \((-R)^n<0\). For \(x=R\), we obtain

\[
\begin{aligned}
P(R)
&\geq a_nR^n-\left|P(R)-a_nR^n\right|\\
&>a_nR^n-\frac12|a_nR^n|
=\frac12a_nR^n>0.
\end{aligned}
\]

For \(x=-R\), we obtain

\[
\begin{aligned}
P(-R)
&\leq a_n(-R)^n+\left|P(-R)-a_n(-R)^n\right|\\
&<a_n(-R)^n+\frac12|a_n(-R)^n|
=\frac12a_n(-R)^n<0.
\end{aligned}
\]

Thus,

\[
P(-R)<0<P(R).
\]

Because \(P\) is continuous on \([-R,R]\), the Intermediate Value Theorem gives some \(c\in[-R,R]\) such that

\[
P(c)=0.
\]

Therefore, \(P\) has at least one real root. \(\square\)

\subsubsection{Proof of the Intermediate Value Theorem}\label{proof-of-the-intermediate-value-theorem}

Let \(f\) be continuous on \([a,b]\), and suppose first that

\[
f(a)<y_0<f(b).
\]

Define

\[
S=\{x\in[a,b]:f(x)\leq y_0\}.
\]

The set \(S\) is nonempty because \(a\in S\), and it is bounded above by \(b\). By the completeness of \(\mathbb{R}\), the supremum

\[
c=\sup S
\]

exists. Since \(f(a)<y_0<f(b)\), we have \(a<c<b\).

We claim that \(f(c)=y_0\). Suppose first that \(f(c)<y_0\). By continuity of \(f\) at \(c\), there exists \(\delta>0\) such that

\[
|x-c|<\delta\quad\Longrightarrow\quad f(x)<y_0.
\]

Choosing \(x>c\) sufficiently close to \(c\) gives \(x\in S\), contradicting the fact that \(c\) is an upper bound of \(S\).

On the other hand, suppose that \(f(c)>y_0\). By continuity of \(f\) at \(c\), there exists \(\delta>0\) such that

\[
|x-c|<\delta\quad\Longrightarrow\quad f(x)>y_0.
\]

Since \(c=\sup S\), there is some \(x\in S\) with \(c-\delta<x\leq c\). This contradicts \(f(x)>y_0\).

Therefore neither \(f(c)<y_0\) nor \(f(c)>y_0\) is possible, and hence

\[
f(c)=y_0.
\]

If \(y_0=f(a)\) or \(y_0=f(b)\), the conclusion is immediate by taking \(c=a\) or \(c=b\). If \(f(b)<f(a)\), apply the same argument to \(-f\). Therefore, for every value \(y_0\) between \(f(a)\) and \(f(b)\), there exists \(c\in[a,b]\) such that

\[
f(c)=y_0.
\]

\subsection{Lecture 6}\label{lecture-6}

\subparagraph{2026.9.16}\label{section-2}

\subsection{Limits Involving Infinity}\label{limits-involving-infinity}

\subsubsection{\texorpdfstring{Finite limits as \(x\to+\infty\)}{Finite limits as x\textbackslash to+\textbackslash infty}}\label{finite-limits-as-xtoinfty}

We say that

\[
\lim_{x\to+\infty}f(x)=L
\]

if, for every \(\varepsilon>0\), there exists \(M>0\) such that

\[
x>M\quad\Longrightarrow\quad |f(x)-L|<\varepsilon.
\]

In words, \(f(x)\) becomes arbitrarily close to \(L\) when \(x\) is sufficiently large.

Similarly,

\[
\lim_{x\to-\infty}f(x)=L
\]

if, for every \(\varepsilon>0\), there exists \(M<0\) such that

\[
x<M\quad\Longrightarrow\quad |f(x)-L|<\varepsilon.
\]

\subsubsection{Examples}\label{examples}

For the reciprocal function,

\[
\lim_{x\to+\infty}\frac{1}{x}=0,
\qquad
\lim_{x\to-\infty}\frac{1}{x}=0.
\]

Indeed, if \(|x|>1/\varepsilon\), then

\[
\left|\frac{1}{x}-0\right|=\frac{1}{|x|}<\varepsilon.
\]

For a rational function whose numerator has lower degree than its denominator,

\[
\lim_{x\to\pm\infty}\frac{P(x)}{Q(x)}=0,
\qquad \deg P<\deg Q.
\]

\subsubsection{\texorpdfstring{Infinite limits as \(x\to c\)}{Infinite limits as x\textbackslash to c}}\label{infinite-limits-as-xto-c}

We write

\[
\lim_{x\to c}f(x)=+\infty
\]

if, for every \(A>0\), there exists \(\delta>0\) such that

\[
0<|x-c|<\delta
\quad\Longrightarrow\quad
f(x)>A.
\]

Similarly,

\[
\lim_{x\to c}f(x)=-\infty
\]

if, for every \(B>0\), there exists \(\delta>0\) such that

\[
0<|x-c|<\delta
\quad\Longrightarrow\quad
f(x)<-B.
\]

For example,

\[
\lim_{x\to0}\frac{1}{x^2}=+\infty,
\]

because \(1/x^2\) can be made larger than any prescribed positive number by choosing \(x\) sufficiently close to \(0\), with \(x\neq0\).

\subsubsection{\texorpdfstring{Infinite limits as \(x\to\pm\infty\)}{Infinite limits as x\textbackslash to\textbackslash pm\textbackslash infty}}\label{infinite-limits-as-xtopminfty}

The notation

\[
\lim_{x\to+\infty}f(x)=+\infty
\]

means that for every \(A>0\), there exists \(M>0\) such that

\[
x>M\quad\Longrightarrow\quad f(x)>A.
\]

The other cases are defined analogously. For instance, if \(n\) is odd and \(a_n>0\), then

\[
\lim_{x\to+\infty}a_nx^n=+\infty,
\qquad
\lim_{x\to-\infty}a_nx^n=-\infty.
\]

\subsection{Exercise 1.1.30}\label{exercise-1.1.30}

\subsubsection{(a)}\label{a}

Prove, using the \(N\)-\(\delta\) language, that

\[
\lim_{x\to0^+}\frac{1}{\sqrt[3]{x}}=+\infty.
\]

\subsubsection{(b)}\label{b}

Find the limit

\[
\lim_{x\to+\infty}\left(\sqrt{x+1}-\sqrt{x-1}\right),
\]

and prove your result using the \(\varepsilon\)-\(N\) language.

\subsubsection{Solution}\label{solution}

\paragraph{(a)}\label{a-1}

Let \(N>0\) be arbitrary. Choose

\[
\delta=\frac{1}{N^3}.
\]

If \(0<x<\delta\), then

\[
\sqrt[3]{x}<\sqrt[3]{\delta}=\frac{1}{N}.
\]

Therefore,

\[
\frac{1}{\sqrt[3]{x}}>\frac{1}{\sqrt[3]{\delta}}=N.
\]

Hence, for every \(N>0\), there exists \(\delta>0\) such that

\[
0<x<\delta
\quad\Longrightarrow\quad
\frac{1}{\sqrt[3]{x}}>N.
\]

Thus,

\[
\boxed{\displaystyle\lim_{x\to0^+}\frac{1}{\sqrt[3]{x}}=+\infty}.
\]

\paragraph{(b)}\label{b-1}

Let

\[
f(x)=\sqrt{x+1}-\sqrt{x-1}.
\]

For \(x\geq1\), rationalizing gives

\[
f(x)
=\frac{(x+1)-(x-1)}{\sqrt{x+1}+\sqrt{x-1}}
=\frac{2}{\sqrt{x+1}+\sqrt{x-1}}.
\]

Since \(\sqrt{x+1}+\sqrt{x-1}\to+\infty\) as \(x\to+\infty\), we obtain

\[
\lim_{x\to+\infty}f(x)=0.
\]

To prove this using the \(\varepsilon\)-\(N\) language, let \(\varepsilon>0\) be given. Choose

\[
N=1+\frac{1}{\varepsilon^2}.
\]

If \(x>N\), then \(x>1\) and

\[
0<f(x)
=\frac{2}{\sqrt{x+1}+\sqrt{x-1}}
\leq\frac{1}{\sqrt{x-1}}.
\]

Moreover, \(x>N\) implies \(x-1>1/\varepsilon^2\), so

\[
\left|f(x)-0\right|=f(x)
\leq\frac{1}{\sqrt{x-1}}<\varepsilon.
\]

Therefore, for every \(\varepsilon>0\), there exists \(N>0\) such that

\[
x>N\quad\Longrightarrow\quad
\left|\sqrt{x+1}-\sqrt{x-1}\right|<\varepsilon.
\]

Hence,

\[
\boxed{\displaystyle
\lim_{x\to+\infty}\left(\sqrt{x+1}-\sqrt{x-1}\right)=0}.
\]

\subsection{Exercise 1.1.28}\label{exercise-1.1.28}

All the Limit Laws remain true when we replace

\[
\lim_{x\to c}
\]

by

\[
\lim_{x\to+\infty}
\quad\text{or}\quad
\lim_{x\to-\infty}.
\]

For example, if

\[
\lim_{x\to+\infty}f(x)=L,
\qquad
\lim_{x\to+\infty}g(x)=M,
\]

then

\[
\lim_{x\to+\infty}(f(x)+g(x))=L+M,
\]

and, provided \(M\neq0\),

\[
\lim_{x\to+\infty}\frac{f(x)}{g(x)}=\frac{L}{M}.
\]

The same statements hold as \(x\to-\infty\).

\subsection{Asymptotes}\label{asymptotes}

The graph of \(y=f(x)\) on an interval \(I\) is

\[
\{(x,y):x\in I,\ y=f(x)\}.
\]

A line is called an \textbf{asymptote} of the graph if the distance between the graph and the line approaches \(0\) as the point on the graph moves to infinity or approaches a finite endpoint.

\subsubsection{Horizontal asymptotes}\label{horizontal-asymptotes}

A line \(y=b\) is a horizontal asymptote of \(y=f(x)\) as \(x\to+\infty\) if

\[
\lim_{x\to+\infty}f(x)=b.
\]

Similarly, \(y=b\) is a horizontal asymptote as \(x\to-\infty\) if

\[
\lim_{x\to-\infty}f(x)=b.
\]

\paragraph{Example 1.1.32}\label{example-1.1.32}

Find the horizontal asymptotes of

\[
f(x)=\frac{\sin x+\cos x}{1+x^2},
\]

and

\[
g(x)=\sqrt{1-\frac{1}{x^2}},\qquad |x|\geq1.
\]

Since \(|\sin x+\cos x|\leq2\),

\[
\left|\frac{\sin x+\cos x}{1+x^2}\right|
\leq\frac{2}{1+x^2}\longrightarrow0.
\]

Thus, \(y=0\) is a horizontal asymptote of \(f\) as \(x\to\pm\infty\).

For \(g\),

\[
\lim_{x\to\pm\infty}g(x)
=\sqrt{1-\lim_{x\to\pm\infty}\frac1{x^2}}
=1.
\]

Thus, \(y=1\) is a horizontal asymptote of \(g\) as \(x\to\pm\infty\).

\subsubsection{Sketches of horizontal asymptotes}\label{sketches-of-horizontal-asymptotes}

\begin{figure}[htbp]
\centering
\includegraphics[width=0.9\linewidth,height=0.32\textheight,keepaspectratio]{image-13.png}
\end{figure}

\begin{figure}[htbp]
\centering
\includegraphics[width=0.9\linewidth,height=0.32\textheight,keepaspectratio]{image-14.png}
\end{figure}

\subsubsection{Vertical asymptotes}\label{vertical-asymptotes}

A line \(x=a\) is a vertical asymptote of \(y=f(x)\) as \(x\to a^+\) if

\[
\lim_{x\to a^+}f(x)=+\infty
\quad\text{or}\quad
\lim_{x\to a^+}f(x)=-\infty.
\]

The analogous definition applies as \(x\to a^-\). In the notation \(\pm\infty\), the sign may be either \(+\) or \(-\).

\paragraph{Example 1.1.33}\label{example-1.1.33}

Consider

\[
h(x)=\sqrt{1-\frac{\sin x}{x^2}},
\]

on its domain

\[
(-\infty,0)\cup[a,+\infty),
\]

where the boundary point \(a>0\) is determined by

\[
a^2=\sin a,
\]

where \(0<a<1\). As \(x\to0^-\), we have \(\sin x\sim x\), and therefore

\[
1-\frac{\sin x}{x^2}
\sim 1-\frac1x\longrightarrow+\infty.
\]

Consequently,

\[
\lim_{x\to0^-}h(x)=+\infty,
\]

so \(x=0\) is a vertical asymptote. Near \(x=a\) from the right, the expression under the square root approaches \(0\), so \(x=a\) is not a vertical asymptote.

This example emphasizes that a domain endpoint is not automatically a vertical asymptote: the function itself must tend to \(\pm\infty\).

\subsubsection{Sketches of vertical asymptotes}\label{sketches-of-vertical-asymptotes}

\begin{figure}[htbp]
\centering
\includegraphics[width=0.9\linewidth,height=0.32\textheight,keepaspectratio]{image-15.png}
\end{figure}

\begin{figure}[htbp]
\centering
\includegraphics[width=0.9\linewidth,height=0.32\textheight,keepaspectratio]{image-16.png}
\end{figure}
\subsubsection{Oblique asymptotes}

A line

\[
y=ax+b,\qquad a\neq0,
\]

is an oblique asymptote of \(y=f(x)\) as \(x\to+\infty\) if

\[
\lim_{x\to+\infty}[f(x)-ax-b]=0.
\]

The definition as \(x\to-\infty\) is analogous.

When the limits exist, the coefficients can be found from

\[
a=\lim_{x\to\pm\infty}\frac{f(x)}{x},
\qquad
b=\lim_{x\to\pm\infty}[f(x)-ax].
\]

\paragraph{Example 1.1.34}\label{example-1.1.34}

Find the oblique asymptote of

\[
f(x)=\frac{x^2-3}{2x-4}.
\]

Polynomial division gives

\[
\begin{aligned}
f(x)
&=\frac{x^2-4+1}{2(x-2)}\\
&=\frac12(x+2)+\frac{1}{2(x-2)}.
\end{aligned}
\]

Therefore,

\[
f(x)-\left(\frac{x}{2}+1\right)
=\frac{1}{2(x-2)}\longrightarrow0
\qquad(x\to\pm\infty).
\]

The oblique asymptote is

\[
y=\frac{x}{2}+1.
\]

Also, \(x=2\) is a vertical asymptote.

\begin{figure}[htbp]
\centering
\includegraphics[width=0.9\linewidth,height=0.32\textheight,keepaspectratio]{image-17.png}
\end{figure}

\paragraph{Example 1.1.35}\label{example-1.1.35}

Let

\[
f(x)=\frac{x^3}{(x+2)(x-1)}.
\]

The vertical asymptotes are

\[
x=-2,\qquad x=1.
\]

For the oblique asymptote,

\[
\frac{f(x)}{x}
=\frac{x^2}{(x+2)(x-1)}\longrightarrow1,
\]

and

\[
f(x)-x
=\frac{-x^2+2x}{(x+2)(x-1)}\longrightarrow-1.
\]

Hence the oblique asymptote is

\[
y=x-1
\]

as \(x\to\pm\infty\).

\begin{figure}[htbp]
\centering
\includegraphics[width=0.9\linewidth,height=0.32\textheight,keepaspectratio]{image-18.png}
\end{figure}

\subsection{Exercise 1.1.36}\label{exercise-1.1.36}

Let

\[
f(x)=\frac1x\sin\frac1x.
\]

Does the graph of \(f\) have a vertical asymptote at \(x=0\)? Explain your answer using the definition of a vertical asymptote.

Since

\[
\left|\frac1x\sin\frac1x\right|\leq\frac1{|x|},
\]

this estimate alone does not determine an infinite limit. To see the oscillation explicitly, define

\[
x_n=\frac{1}{\frac{\pi}{2}+2\pi n},
\qquad
y_n=\frac{1}{\frac{3\pi}{2}+2\pi n}.
\]

Then \(x_n,y_n\to0^+\) and

\[
f(x_n)=\frac{\pi}{2}+2\pi n\longrightarrow+\infty,
\]

while

\[
f(y_n)=-\left(\frac{3\pi}{2}+2\pi n\right)\longrightarrow-\infty.
\]

Thus, as \(x\to0^+\), the function tends neither to \(+\infty\) nor to \(-\infty\). Applying the same argument to \(-x_n\) and \(-y_n\) gives the same conclusion as \(x\to0^-\). Therefore, \(x=0\) is not a vertical asymptote in the usual one-sided sense.

\begin{figure}[htbp]
\centering
\includegraphics[width=0.9\linewidth,height=0.32\textheight,keepaspectratio]{image-19.png}
\end{figure}

\subsection{More Exercises}\label{more-exercises}

Find the limits and use the \(\varepsilon\)-\(N\) language to prove your results.

\subsubsection{(a)}\label{a-2}

\[
\lim_{x\to+\infty}\sqrt{x}\left(\sqrt{x+1}-\sqrt{x-1}\right).
\]

\subsubsection{(b)}\label{b-2}

\[
\lim_{x\to+\infty}
\sqrt{x+\sqrt{x+\frac{\sqrt{x}}{\sqrt{x+1}}}}.
\]

\section{Lecture 7}\label{lecture-7}

\subparagraph{2026 9.17}\label{section-3}

\subsection{1.7 Sequences and limits}\label{sequences-and-limits}

\subsubsection{Definition 1.1.12 (Limit of a sequence)}\label{definition-1.1.12-limit-of-a-sequence}

A sequence \(\{a_n\}\) converges to \(L\) if for every \(\varepsilon>0\) there exists an integer \(N\) such that

\[
 n>N \quad\Longrightarrow\quad |a_n-L|<\varepsilon.
\]

If no such number \(L\) exists, we say that the sequence diverges.

We write

\[
\lim_{n\to\infty}a_n=L,
\qquad\text{or}\qquad a_n\to L.
\]

\subsubsection{Exercise 1.1.39}\label{exercise-1.1.39}

Show that if a sequence converges, then the limit is unique.

\paragraph{Proof}\label{proof-1}

Suppose that

\[
 a_n\to L_1,
\qquad a_n\to L_2.
\]

Let \(\varepsilon>0\) be arbitrary. Then there exist \(N_1,N_2\) such that

\[
 n>N_1 \Rightarrow |a_n-L_1|<\varepsilon/2,
\]

and

\[
 n>N_2 \Rightarrow |a_n-L_2|<\varepsilon/2.
\]

For \(n>\max\{N_1,N_2\}\),

\[
 |L_1-L_2|
\le |a_n-L_1|+|a_n-L_2|
< \varepsilon.
\]

Since this holds for every \(\varepsilon>0\), we must have \(L_1=L_2\).

Hence the limit is unique.

\subsubsection{Example 1.1.40}\label{example-1.1.40}

Investigate the convergence or divergence of the following sequences.

\paragraph{\texorpdfstring{(a) \((-1)^n\)}{(a) (-1)\^{}n}}\label{a--1n}

This sequence alternates between \(-1\) and \(1\), so it does not approach any single number. Therefore,

\[
(-1)^n\text{ diverges.}
\]

\paragraph{\texorpdfstring{(b) \(\displaystyle \frac{\sin n-n^2}{n^3+1}\)}{(b) \textbackslash displaystyle \textbackslash frac\{\textbackslash sin n-n\^{}2\}\{n\^{}3+1\}}}\label{b-displaystyle-fracsin-n-n2n31}

Since

\[
\frac{\sin n-n^2}{n^3+1}
=
\frac{\sin n}{n^3+1}-\frac{n^2}{n^3+1},
\]

and both terms go to \(0\), we have

\[
\lim_{n\to\infty}\frac{\sin n-n^2}{n^3+1}=0.
\]

\paragraph{\texorpdfstring{(c) \(\displaystyle \frac{(-1)^{[\sqrt n]}}{n}\)}{(c) \textbackslash displaystyle \textbackslash frac\{(-1)\^{}\{{[}\textbackslash sqrt n{]}\}\}\{n\}}}\label{c-displaystyle-frac-1sqrt-nn}

Here \([\sqrt n]\) is the integer part of \(\sqrt n\), and the numerator is always \(\pm1\). Since dividing by \(n\) makes the magnitude go to \(0\),

\[
\lim_{n\to\infty}\frac{(-1)^{[\sqrt n]}}{n}=0.
\]

\subsection{Bounded sequences}\label{bounded-sequences}

\subsubsection{Definition}\label{definition}

A sequence \(\{a_n\}\) is called:

\begin{itemize}
\tightlist
\item
  bounded above if there exists a constant \(M\) such that \(a_n\le M\) for all \(n\);
\item
  bounded below if there exists a constant \(m\) such that \(a_n\ge m\) for all \(n\);
\item
  bounded if there exists a constant \(C>0\) such that \(|a_n|\le C\) for all \(n\).
\end{itemize}

Equivalently, \(\{a_n\}\) is bounded iff there exist \(m,M\) with

\[
 m\le a_n\le M \qquad \forall n\ge1.
\]

\subsubsection{Theorem}\label{theorem}

Every convergent sequence is bounded.

\paragraph{Proof}\label{proof-2}

If \(a_n\to L\), then for \(\varepsilon=1\) there exists \(N\) such that

\[
 n>N \Rightarrow |a_n-L|<1.
\]

Thus for all \(n>N\),

\[
 |a_n|\le |L|+1.
\]

Let

\[
 C=\max\{|a_1|,\dots,|a_N|,|L|+1\}.
\]

Then \(|a_n|\le C\) for all \(n\), so the sequence is bounded.

\subsection{Divergence to infinity}\label{divergence-to-infinity}

\subsubsection{Definition}\label{definition-1}

We say that \(a_n\to +\infty\) if for every real number \(M\) there exists an integer \(N\) such that

\[
 n\ge N \Rightarrow a_n>M.
\]

We say that \(a_n\to -\infty\) if for every real number \(m>0\) there exists an integer \(N\) such that

\[
 n\ge N\Rightarrow a_n<-m.
\]

\subsubsection{Example 1.1.41}\label{example-1.1.41}

Investigate the behavior of the following sequences as \(n\to\infty\).

\paragraph{\texorpdfstring{(a) \((-1)^n n\)}{(a) (-1)\^{}n n}}\label{a--1n-n}

This sequence alternates sign and grows in magnitude, so it is unbounded and does not converge to a finite limit. It also does not tend to \(+\infty\) or \(-\infty\) in the usual sense because the signs alternate.

\paragraph{\texorpdfstring{(b) \(a^n/n!\) with \(a>0\)}{(b) a\^{}n/n! with a\textgreater0}}\label{b-ann-with-a0}

For fixed \(a>0\), the factorial grows faster than any exponential, so

\[
\frac{a^n}{n!}\to0.
\]

\paragraph{\texorpdfstring{(c) \(n^n/n!\)}{(c) n\^{}n/n!}}\label{c-nnn}

Since

\[
\frac{n^n}{n!}=\prod_{k=1}^n\frac{n}{k},
\]

this sequence grows very rapidly, so

\[
\frac{n^n}{n!}\to +\infty.
\]

\subsection{Limit laws for sequences}\label{limit-laws-for-sequences}

If

\[
\lim_{n\to\infty}a_n=A,
\qquad
\lim_{n\to\infty}b_n=B,
\]

then the following hold:

\begin{enumerate}
\def\labelenumi{\arabic{enumi}.}
\tightlist
\item
  \(\displaystyle \lim_{n\to\infty}(a_n+b_n)=A+B\);
\item
  \(\displaystyle \lim_{n\to\infty}(a_n-b_n)=A-B\);
\item
  \(\displaystyle \lim_{n\to\infty}(k a_n)=kA\) for any constant \(k\);
\item
  \(\displaystyle \lim_{n\to\infty}(a_n b_n)=AB\);
\item
  if \(B\neq0\), then \[
  \lim_{n\to\infty}\frac{a_n}{b_n}=\frac{A}{B}.
  \]
\end{enumerate}

\subsection{Exercise 1.1.43}\label{exercise-1.1.43}

Let \(\{a_n\}\) and \(\{b_n\}\) be convergent sequences and assume \(a_n\ge b_n\) for all sufficiently large \(n\). Prove that

\[
\lim_{n\to\infty}a_n\ge \lim_{n\to\infty}b_n.
\]

\subsubsection{Proof}\label{proof-3}

Let

\[
 A=\lim_{n\to\infty}a_n,
\qquad B=\lim_{n\to\infty}b_n.
\]

Suppose, for contradiction, that \(A<B\). Set

\[
 \varepsilon=\frac{B-A}{2}>0.
\]

Then there exists \(N\) such that for all \(n>N\),

\[
 |a_n-A|<\varepsilon,
\qquad |b_n-B|<\varepsilon.
\]

Hence for \(n>N\),

\[
 a_n<A+\varepsilon=\frac{A+B}{2},
\]

and

\[
 b_n>B-\varepsilon=\frac{A+B}{2}.
\]

So \(a_n<b_n\) for all sufficiently large \(n\), contradicting the assumption that \(a_n\ge b_n\) eventually. Therefore,

\[
 A\ge B.
\]

\subsection{Exercise 1.1.43 (continued)}\label{exercise-1.1.43-continued}

Assume \(a_n\le b_n\le c_n\) for all \(n\ge N_0\) and

\[
 \lim_{n\to\infty}a_n = \lim_{n\to\infty}c_n = L.
\]

Then

\[
\lim_{n\to\infty}b_n = L.
\]

\subsubsection{Proof}\label{proof-4}

Let \(\varepsilon>0\) be given. Since \(a_n\to L\) and \(c_n\to L\), there exist \(N_1,N_2\) such that

\[
 n>N_1 \Rightarrow |a_n-L|<\varepsilon,
\]

and

\[
 n>N_2 \Rightarrow |c_n-L|<\varepsilon.
\]

Let \(N=\max\{N_0,N_1,N_2\}\). Then for all \(n>N\),

\[
 L-\varepsilon<a_n\le b_n\le c_n<L+\varepsilon.
\]

Hence

\[
 |b_n-L|<\varepsilon,
\]

for all \(n>N\). Therefore, \(b_n\to L\).

This is the sequence version of the Sandwich Theorem.

\subsection{Exercise 1.1.44}\label{exercise-1.1.44}

Let \(\{a_n\}\) be a sequence with \(a_n\to L\), and let \(f\) be a function continuous at \(L\) and defined at all \(a_n\). Then

\[
 f(a_n)\to f(L).
\]

\subsubsection{Proof}\label{proof-5}

Let \(\varepsilon>0\) be given. Since \(f\) is continuous at \(L\), there exists \(\delta>0\) such that

\[
 |x-L|<\delta \Rightarrow |f(x)-f(L)|<\varepsilon.
\]

Because \(a_n\to L\), there exists \(N\) such that

\[
 n>N \Rightarrow |a_n-L|<\delta.
\]

Therefore, for all \(n>N\),

\[
 |f(a_n)-f(L)|<\varepsilon.
\]

So \(f(a_n)\to f(L)\).

\subsection{Proposition 1.1.13 (function limits and sequence limits)}\label{proposition-1.1.13-function-limits-and-sequence-limits}

Let \(I\) be an interval containing \(x_0\) in its interior and let \(f\) be defined on \(I\setminus\{x_0\}\). Then

\[
\lim_{x\to x_0}f(x)=A
\]

exists if and only if for every sequence \(\{a_n\}\subset I\setminus\{x_0\}\) with \(a_n\to x_0\), the limit

\[
\lim_{n\to\infty}f(a_n)
\]

exists and equals \(A\).

\subsubsection{Proof sketch}\label{proof-sketch}

\begin{itemize}
\tightlist
\item
  If \(f(x)\to A\) as \(x\to x_0\), then for any sequence \(a_n\to x_0\), the values \(f(a_n)\) must also approach \(A\).
\item
  Conversely, if every sequence approaching \(x_0\) produces the same limiting value \(A\), then \(f(x)\) cannot fail to approach \(A\); otherwise one could construct a sequence \(x_n\to x_0\) with \(|f(x_n)-A|\ge\varepsilon_0\) for all \(n\), contradicting the assumption.
\end{itemize}

Thus, the limit of a function can be tested using sequences.

\subsection{Exercise 1.1.45}\label{exercise-1.1.45}

Let \(f\) be a function defined on \([M,\infty)\). Prove:

\subsubsection{(a)}\label{a-3}

\[
\lim_{x\to\infty}f(x)=L
\]

if and only if for every sequence \(a_n\to\infty\),

\[
\lim_{n\to\infty}f(a_n)=L.
\]

\subsubsection{(b)}\label{b-3}

\[
\lim_{x\to\infty}f(x)=+\infty
\]

if and only if for every sequence \(a_n\to\infty\),

\[
\lim_{n\to\infty}f(a_n)=+\infty.
\]

These results show that the behavior of a function at infinity can be studied through sequences tending to infinity.

\subsection{1.8 Least upper bound and greatest lower bound}\label{least-upper-bound-and-greatest-lower-bound}

\subsubsection{Definition}\label{definition-2}

Let \(\{a_n\}\) be a sequence bounded above. A number \(L\) is called the least upper bound, or supremum, of \(\{a_n\}\) if:

\begin{enumerate}
\def\labelenumi{\arabic{enumi}.}
\tightlist
\item
  \(a_n\le L\) for all \(n\);
\item
  if \(M\) is any upper bound, then \(L\le M\).
\end{enumerate}

We write

\[
 L=\sup_{n\ge1}a_n.
\]

Similarly, if \(\{a_n\}\) is bounded below, a number \(l\) is called the greatest lower bound, or infimum, if:

\begin{enumerate}
\def\labelenumi{\arabic{enumi}.}
\tightlist
\item
  \(a_n\ge l\) for all \(n\);
\item
  if \(m\) is any lower bound, then \(l\ge m\).
\end{enumerate}

We write

\[
 l=\inf_{n\ge1}a_n.
\]

\subsubsection{Remarks}\label{remarks}

\begin{itemize}
\tightlist
\item
  For the supremum: for every \(\varepsilon>0\), there exists \(n\) such that \[
  a_n>L-\varepsilon.
  \]
\item
  For the infimum: for every \(\varepsilon>0\), there exists \(n\) such that \[
  a_n<l+\varepsilon.
  \]
\end{itemize}

\subsubsection{Theorem 1.1.15}\label{theorem-1.1.15}

Any sequence bounded above has a supremum, and any sequence bounded below has an infimum.

This is the completeness property of the real numbers.

\subsection{Example 1.1.47}\label{example-1.1.47}

Find the infimum and supremum of the following sequences.

\subsubsection{\texorpdfstring{(a) \(\left\{(-1)^n/n\right\}_{n\ge1}\)}{(a) \textbackslash left\textbackslash\{(-1)\^{}n/n\textbackslash right\textbackslash\}\_\{n\textbackslash ge1\}}}\label{a-left-1nnright_nge1}

\[
\inf=-\frac12,\qquad \sup=\frac12.
\]

\subsubsection{\texorpdfstring{(b) \(\{(-1)^n n\}_{n\ge1}\)}{(b) \textbackslash\{(-1)\^{}n n\textbackslash\}\_\{n\textbackslash ge1\}}}\label{b--1n-n_nge1}

This sequence is unbounded above and below, so

\[
\sup=+\infty,
\qquad\inf=-\infty.
\]

\subsubsection{\texorpdfstring{(c) \(\{n(-1)^n\}_{n\ge1}\)}{(c) \textbackslash\{n(-1)\^{}n\textbackslash\}\_\{n\textbackslash ge1\}}}\label{c-n-1n_nge1}

Again unbounded, so

\[
\sup=+\infty,
\qquad\inf=-\infty.
\]

\subsubsection{\texorpdfstring{(d) \(a_1=\sqrt2\), \(a_{n+1}=\sqrt{2+a_n}\)}{(d) a\_1=\textbackslash sqrt2, a\_\{n+1\}=\textbackslash sqrt\{2+a\_n\}}}\label{d-a_1sqrt2-a_n1sqrt2a_n}

This sequence is increasing and bounded above by \(2\), so it converges to the positive solution of

\[
 x=\sqrt{2+x},
\]

namely

\[
 x^2-x-2=0 \Rightarrow x=2.
\]

Thus,

\[
\sup=2,
\qquad\inf=\sqrt2.
\]

\section{Lecture 8}\label{lecture-8}

\subparagraph{2026 9.17}\label{section-4}

\end{document}
